\documentclass[12pt]{article}
\usepackage[utf8]{inputenc}
\usepackage[T1]{fontenc}
\usepackage[margin=1in]{geometry}
\usepackage{setspace}
\usepackage{xcolor}
\usepackage[colorlinks=true, urlcolor=blue!60!black, linkcolor=blue!60!black]{hyperref}
\setlength{\parindent}{0em}
\setlength{\parskip}{0.6em}
\renewcommand{\baselinestretch}{1.0}
\begin{document}
\begin{center}
{\LARGE \textit{Trump v. Mazars USA, LLP}, 591 U.S. 848 (2020)\par}
\vspace{0.5em}
{\large Case Brief --- Two-Page Edition\par}
\end{center}

\section{Memory Jogger}
House subpoenas for a sitting President's personal financial records require a four-factor separation-of-powers balancing, neither executive-privilege rigor nor ordinary deference (\href{https://supreme.justia.com/cases/federal/us/591/19-715/}{Justia}).

\section{Detailed Case Facts}
In 2019 three House committees subpoenaed Deutsche Bank, Capital One, and the accounting firm Mazars USA for about a decade of President Trump's and his businesses' financial records, citing money-laundering, foreign-influence, ethics, and financial-disclosure legislation.

\section{Procedural History}
The President sued personally, without asserting executive privilege. The D.C. and Second Circuits upheld the subpoenas, over dissents by Judges Rao and Livingston.

\section{Judicial Votes}
Vacated and remanded, 7--2. Roberts wrote for Ginsburg, Breyer, Sotomayor, Kagan, Gorsuch, and Kavanaugh; Thomas and Alito dissented.

\section{Holding}
Courts must weigh (1) whether other sources could supply the information, (2) whether the subpoena is no broader than reasonably necessary, (3) the evidence of a valid legislative purpose, and (4) the burdens on the President.

\section{Analysis of Opinions}
\textbf{Majority:} Two centuries of interbranch accommodation deserve weight. \textit{Nixon}'s privilege standard should not be transplanted ``root and branch'' to private papers, but Congress and the President are rival branches, and concerns persist for personal papers held by third parties. \textbf{Thomas:} Congress has no power to subpoena private documents except through impeachment. \textbf{Alito:} such subpoenas are ``inherently suspicious''; the House should have to justify them in detail.

\section{Examples: Future Applications}
\textbf{Same-Side:} a committee subpoenas a decade of the President's medical records for a general disability statute; enforcement should be refused. \textbf{Opposite-Side:} a state grand jury subpoena under \textit{Trump v.\ Vance} needs no heightened showing. \textbf{Fence-Sitter:} a narrow subpoena aimed at a former President's disclosures.

\section{Critique}
The four-factor test offers little guidance and invites delay that outlasts a two-year House subpoena. A county prosecutor in \textit{Vance} faced a lower bar than Congress, and Thomas's view could push Congress toward premature impeachment inquiries.

\section{Key Quotations}
\begin{quote}
``The President is the only person who alone composes a branch of government. As a result, there is not always a clear line between his personal and official affairs.''
\end{quote}
\noindent --- Chief Justice Roberts, Opinion of the Court.

\end{document}
